\itemhead{习题4.1 · 第11题（教材p.105）}{XJ-B1-C04-HW-P105-Y0490}
\label{XJ-B1-C04-HW-P105-Y0490}
{\footnotesize 湘教版教材p.105 / PDF第112页 / 习题4.1 / 第11题}
\field{原题}{已知 $f(x)$ 是奇函数，且当 $x\ge0$ 时，$f(x)=x^{2/5}$，求当 $x<0$ 时 $f(x)$ 的表达式。}
\field{参考答案}{$x<0$ 时 $f(x)=-f(-x)=-(-x)^{2/5}=-|x|^{2/5}$，不能漏负号。}
{\small \field{关联知识}{F08 奇偶性、对称与函数性质应用（教材帮p.113、114、115 / PDF122、123、124）；P03 幂函数的概念、图象与性质（教材帮p.125、126、127 / PDF134、135、136）}
\field{方法与易错点}{检查定义域，再代入 $-x$；求另一侧解析式时同时写出 $-x$ 所在区间。 易错：只看公式不看定义域；用一对点证明奇偶；负半轴代换遗漏最外层负号。}
\field{答案依据}{答案册相关题组：PDF第31页（册内页30）；按当前题库原题逐项重算，不按旧版题号机械对套。}}
